\subsubsection*{3.2 Die Einheitengruppe $\Zm{m}^*$ von $\Zm{m}$}


\begin{bemerkung}
% KORRIGIERT: G war nicht eingeführt
Die Einheiten sind genau die Elemente $k \in \Zm{m}$ mit $\ggT(k, m) = 1$.\\
$(\Zm{m}^*, \cdot\,)$ ist eine abelsche Gruppe. Für eine endliche Gruppe $G$ ist
\[
  \operatorname{ord}(G) = |G|, \qquad\qquad \operatorname{ord}(\Zm{m}^*) = \varphi(m).
\]
\end{bemerkung}

\begin{bemerkung}[Kürzungsregel]
% KORRIGIERT: Quantor für x, y und Begründung fehlten
Für Einheiten $a \in \Zm{m}^*$ und alle $x, y \in \Zm{m}$ gilt (Multiplikation mit $a^{-1}$)
\[
  a \cdot x = a \cdot y \quad \Rightarrow \quad x = y \qquad\qquad \text{Kürzungsregel}
\]
\end{bemerkung}

% KORRIGIERT: Original "Sei p eine Primzahl. Dann ist": Voraussetzung a fehlte; {a1, a2, ...} statt {a·1, a·2, ...};
% "=" statt "\equiv ... (mod p)"; Kürzen von (p-1)! ohne Begründung
Sei $p$ eine Primzahl und $a \in \Zm{p}^*$. Nach der Kürzungsregel sind $a \cdot 1, a \cdot 2, \ldots, a \cdot (p - 1)$ paarweise verschieden. Dann ist (in $\Zm{p}$)
\begin{align*}
  \Rightarrow\quad & \Zm{p}^* = \{1, 2, 3, \ldots, p - 1\} = \{a \cdot 1, a \cdot 2, a \cdot 3, \ldots, a \cdot (p - 1)\} \\
  \intertext{Dann stimmen auch die Produkte überein}
  \Rightarrow\quad & 1 \cdot 2 \cdot 3 \cdot \ldots \cdot (p - 1) \equiv a^{p-1} \cdot 1 \cdot 2 \cdot 3 \cdot \ldots \cdot (p - 1) \pmod{p} \\
  \intertext{Wegen $\ggT((p - 1)!, p) = 1$ ist $(p - 1)!$ in $\Zm{p}$ invertierbar, man darf also kürzen:}
  \Rightarrow\quad & 1 \equiv a^{p-1} \pmod{p}
\end{align*}

\begin{beispiel}
% KORRIGIERT: {a1, a2, ...} statt {a·1, a·2, ...}
$p = 7$, $a = 3$ $\Rightarrow$ $\{a \cdot 1, a \cdot 2, a \cdot 3, \ldots, a \cdot 6\} = \{3, 6, 2, 5, 1, 4\}$ \ (in $\Zm{7}$)
\end{beispiel}

\begin{satz}[Fermat]
% KORRIGIERT: G nicht benannt; Voraussetzung für die Form in Z fehlte
Sei $p$ eine Primzahl und $a \in \Zm{p}^*$, dann gilt \hfill $a^{\operatorname{ord} G} = 1$ \ (mit $G = \Zm{p}^*$).
\begin{align*}
  & a^{p-1} \equiv 1 \pmod{p} \qquad \text{für alle } a \in \Z \text{ mit } p \nmid a \\
  \text{bzw.}\quad & a^{p-1} = 1 \quad (\text{in } \Zm{p}). \qquad \square
\end{align*}
\end{satz}

Im Beweis haben wir genutzt, dass $\Zm{p}^*$ abelsch ist.

\begin{definition}[Bezeichnung]
% KORRIGIERT: G und a waren nicht eingeführt
Sei $G$ eine endliche Gruppe und $a \in G$. Die kleinste Zahl $k > 0$ mit $a^k = 1$ heißt \textbf{Ordnung von $a$} in $G$. Die Menge
\[
  \langle a \rangle := \{\, a, a^2, a^3, \ldots, a^k = 1 \,\} \subseteq G
\]
ist eine zyklische Gruppe der Ordnung $k$.

Ein Element $g \in G$ mit $\operatorname{ord}(g) = |G|$ heißt erzeugendes Element (oder \textbf{primitiv}) von $G$. Für solches $g$ gilt
\[
  G = \{\, g, g^2, \ldots, g^k = 1 \} = \langle g \rangle \qquad (k = |G|)
\]
Besitzt $G$ ein primitives Element, so heißt $G$ zyklisch. Zyklische Gruppen sind abelsch.
\end{definition}

\begin{beispiel}
$\Zm{7}^* = \{1, 2, \ldots, 6\}$
\begin{align*}
  & 2, 2^2, 2^3, \ldots, 2^6 = 2, 4, 1, 2, 4, 1 && \Rightarrow\ \operatorname{ord}(2) = 3 \\
  & 3, 3^2, 3^3, \ldots, 3^6 = 3, 2, -1, -3, -2, 1 && \operatorname{ord}(3) = 6 \\
  \Rightarrow\quad & 3 \text{ ist primitives Element von } \Zm{7}^*.
\end{align*}
\end{beispiel}

\begin{bemerkung}
Fermat war begeistert von seinem Satz. % KORRIGIERT: Tippfehler "zeige"
Er zeigte, dass die Griechen noch nicht alles wußten.

Euler hat den Satz von Fermat etwa 100 Jahre später auf $\Zm{m}^*$ übertragen.\\
$\operatorname{ord} \Zm{m}^* = \boldsymbol{\varphi}(m)$.
\end{bemerkung}

\begin{satz}[Euler]
Formulierung in $\Zm{m}$
\[
  a^{\varphi(m)} = 1 \qquad\qquad \text{für } a \in \Zm{m}^*
\]
Formulierung in $\Z$
\[
  a^{\varphi(m)} \equiv 1 \pmod{m} \qquad \text{für } \ggT(a, m) = 1
\]
\end{satz}

\begin{beispiel}
% KORRIGIERT: "=" und "\neq" statt "\equiv" und "\not\equiv"; Begründung für phi(2·17) = phi(2)·phi(17) fehlte
Fortsetzung: $\varphi(34) = \varphi(2 \cdot 17) = \varphi(2) \cdot \varphi(17) = 1 \cdot 16 = 16$ \ (da $\ggT(2, 17) = 1$)
\[
  \Rightarrow \quad 12^{16} \equiv -16 \not\equiv 1 \pmod{34}
\]
Kein Wunder: $\ggT(12, 34) = 2$ teilt jede Potenz von 12 mod 34. Der Satz von Euler setzt $\ggT(a, m) = 1$ voraus.
\end{beispiel}

\begin{bemerkung}
Falls $\ggT(a, m) = 1$, geht das Potenzieren schneller:
\begin{align*}
% KORRIGIERT: Bereich von r und Verweis auf Euler fehlten; Modulschreibweise vereinheitlicht
  & k = q \cdot \varphi(m) + r,\ 0 \leq r < \varphi(m), \quad \text{also } k \equiv r \pmod{\varphi(m)} \\
  \Rightarrow\quad & a^k = (a^{\varphi(m)})^q \cdot a^r \equiv 1^q \cdot a^r = a^r \pmod{m} \quad (\text{Euler}) \qquad \square
\end{align*}
\end{bemerkung}

\begin{beispiel}
Gesucht: $11^{16} \pmod{34}$
\begin{align*}
  & 11 \\
% KORRIGIERT: Reste 19, 21, -1 und (mod 34) ergänzt
  & 11^2 = 121 = 3 \cdot 34 + 19 \equiv 19 \pmod{34} \\
  & 11^4 \equiv 19^2 = 361 = 10 \cdot 34 + 21 \equiv 21 \pmod{34} \\
  & 11^8 \equiv 21^2 = 441 = 13 \cdot 34 - 1 \equiv -1 \pmod{34} \\
  & 11^{16} \equiv (-1)^2 = 1 \pmod{34}
\end{align*}
\end{beispiel}
